% Part: computability
% Chapter: computability-theory
% Section: russells-paradox

\documentclass[../../../include/open-logic-section]{subfiles}

\begin{document}

\olfileid{cmp}{thy}{rus}
\olsection{রাসেলের কূটাভাসের সঙ্গে তুলনা}

এই অংশের যুক্তিগুলির সঙ্গে রাসেলের কূটাভাসের সাদৃশ্য ও পার্থক্য তুলনা
করা শিক্ষণীয়:
\begin{enumerate}
\item রাসেলের কূটাভাস: ধরা যাক, $S = \Setabs{x}{x \notin x}$। তাহলে
  $S \in S$ যদি এবং কেবল যদি $S \notin S$—যা বিরোধ।
% BN-SRC-146 TARGET-CORRECTION-BEGIN
\par\smallskip\noindent\textnormal{\small\textbf{উৎস-সংশোধন (BN-SRC-146):} হিমায়িত ইংরেজি পাঠে দ্বিশর্তটির ডান দিকে হঠাৎ বড় হাতের~$X$ এসেছে, যদিও নির্মিত সেটটি~$S$ এবং রাসেলের যুক্তির জন্য দরকার $S \in S$ যদি এবং কেবল যদি $S \notin S$। বাংলা পাঠে ডান দিকেও~$S$ রাখা হয়েছে।} % BN-SRC-146 TARGET-CORRECTION-END

  \emph{সিদ্ধান্ত:} এমন কোনো সেট~$S$ নেই। “সব সেটের সেট”-এর অস্তিত্ব
  অনুমান করা সেটতত্ত্বের অন্যান্য স্বতঃসিদ্ধের সঙ্গে অসঙ্গত।

\item রাসেলের কূটাভাসের একটি পরিবর্তিত রূপ: ধরা যাক, $F$ হলো সব অপেক্ষকের
  সেট থেকে $\{ 0, 1 \}$-এ নিচের সংজ্ঞায় নির্ধারিত “অপেক্ষক”:
  \[
  F(f) =
  \begin{cases}
    1 & \text{যদি $f$, $f$-এর সংজ্ঞাক্ষেত্রে থাকে এবং $f(f) = 0$} \\
    0 & \text{অন্যথায়}
  \end{cases}
  \]
  অনুরূপ যুক্তি দেখায় যে $F(F) = 0$ যদি এবং কেবল যদি $F(F) = 1$—যা
  বিরোধ।

  \emph{সিদ্ধান্ত:} $F$ কোনো অপেক্ষক নয়। “সব অপেক্ষকের সেট” কোনো
  অপেক্ষকের সংজ্ঞাক্ষেত্র হওয়ার পক্ষে অতিরিক্ত বড়।

\item কর্ণীকরণ যুক্তি: ধরা যাক, $f_0$, $f_1$, \dots হলো আংশিক গণনসাধ্য
  অপেক্ষকগুলির তালিকায়ন, এবং $G \colon \Nat \to \{ 0, 1 \}$-কে এভাবে
  সংজ্ঞায়িত করা হয়েছে:
  \[
  G(x) =
  \begin{cases}
    1 & \text{যদি $f_x(x)\downarrow = 0$} \\
    0 & \text{অন্যথায়}
  \end{cases}
  \]
  $G$ গণনসাধ্য হলে কোনো $k$-এর জন্য এটি অপেক্ষক $f_k$। কিন্তু তখন
  $G(k) = 1$ যদি এবং কেবল যদি $G(k) = 0$—যা বিরোধ।

  \emph{সিদ্ধান্ত:} $G$ গণনসাধ্য নয়। লক্ষ করুন, সেটতত্ত্বের
  স্বতঃসিদ্ধগুলি অনুসারে $G$ তবু একটি অপেক্ষক; এখানে কোনো কূটাভাস নেই,
  আছে শুধু একটি স্পষ্টীকরণ।
\end{enumerate}

আংশিক অপেক্ষক, গণনসাধ্য অপেক্ষক, আংশিক গণনসাধ্য অপেক্ষক ইত্যাদির কথা
সহজেই বিভ্রান্তিকর হতে পারে। $\Nat$ থেকে $\Nat$-এ সব আংশিক অপেক্ষকের সেট
বস্তুগুলির একটি বড় সংগ্রহ। তাদের কিছু সর্বত্র-সংজ্ঞায়িত, কিছু গণনসাধ্য,
কিছু সর্বত্র-সংজ্ঞায়িত ও গণনসাধ্য দুটিই, আবার কিছু কোনোটিই নয়। মনে রাখুন,
আমরা আলাদা করে কিছু না বললে “অপেক্ষক” বলতে সর্বত্র-সংজ্ঞায়িত অপেক্ষক
বোঝাই। অতএব আমাদের আছে:
\begin{enumerate}
\item গণনসাধ্য অপেক্ষক
\item সর্বত্র-সংজ্ঞায়িত নয় এমন আংশিক গণনসাধ্য অপেক্ষক
\item গণনসাধ্য নয় এমন অপেক্ষক
\item সর্বত্র-সংজ্ঞায়িতও নয়, গণনসাধ্যও নয় এমন আংশিক অপেক্ষক
\end{enumerate}
বিষয়টি গুছিয়ে নিতে $\Nat$ থেকে $\Nat$-এ সব আংশিক অপেক্ষক বোঝানো একটি
বড় বর্গক্ষেত্র আঁকা যেতে পারে; তারপর যথাক্রমে সর্বত্র-সংজ্ঞায়িত অপেক্ষক
ও গণনসাধ্য আংশিক অপেক্ষক বোঝানো দুটি পরস্পরচ্ছেদী অঞ্চল চিহ্নিত করুন।
চিত্রে এভাবে তৈরি প্রতিটি অঞ্চলের একটি করে বস্তুর বর্ণনা দিতে পারেন কি
না দেখা একটি ভালো অনুশীলনী।

\end{document}
